Standard geographic boundary-length edge weights treat all tract boundaries
equivalently, regardless of the communities on either side.
This paper introduces \emph{economic character} edge weights: adjacent census
tracts with similar economic profiles---both predominantly residential, both
dominated by commercial activity, both industrial---receive heavy edges that
METIS avoids cutting, while tracts of different economic character receive
light edges that invite cuts at natural economic zone transitions.

Economic character is operationalized through the Census Bureau's
Longitudinal Employer-Household Dynamics (LEHD) Origin-Destination Employment
Statistics (LODES) Workplace Area Characteristics (WAC) data.
Three signals are computed per tract from raw NAICS-sector job counts:
\emph{commercial intensity}, \emph{industrial fraction}, and
\emph{jobs per resident}.
Edge similarity is computed as the cosine of the angle between the two tracts'
three-dimensional character vectors, and blended with the geographic boundary
weight at a tunable $\alpha$ parameter (default 0.5).

We evaluate the method on North Carolina ($k=14$), Wisconsin ($k=8$), and
Texas ($k=38$) against the standard geographic weight baseline.
In North Carolina, economic character weights shift the proportionality gap
from $-20.7$\,pp (geographic) to $-6.4$\,pp (economic character), a
14.3\,pp improvement, as the Research Triangle Park and Charlotte commercial
corridor emerge as natural cut seams.
In Wisconsin, where commercial and industrial zones are diffuse, the method
produces no effect.
In Texas, industrial clusters in Houston and Dallas cut against existing
geographic sorting, yielding a slight 2.6\,pp worsening.
The key finding is state-specificity: economic character weights are most
effective when a state's economic geography aligns with---and thus can help
correct---its partisan geographic sorting.
The weight computation is $O(|E|) = O(n)$ for planar tract adjacency graphs
after LODES aggregation and adds under 50\,ms overhead on a 2000-tract state.
